% --- Macro pour une surrogate policy + son ensemble ---
\newcommand{\surrogatepolicy}[7]{%
    % #1 : xshift
    % #2 : index (1,2,3,…)
    % #3 : liste des coordonnées X des points
    % #4 : liste des coordonnées Y des points
    % #5 : liste des coordonnées X des circlepattern
    % #6 : liste des coordonnées Y des circlepattern
    % #7 : liste des diamètres des circlepattern
    \node[policybox] (policy#2) at (#1,0)
        {\scalebox{2}{$\textcolor{myred}{\pi}^{\textcolor{myred}{'}}_{\textcolor{myred}{#2}}$}};

    \begin{scope}[xshift=#1, yshift=-4cm, scale=0.8, transform shape]
        \node[pointset] (box#2) at (0,0) {};
        \begin{scope}
            \clip (box#2.south west) rectangle (box#2.north east);
            % Circlepattern avec diamètres spécifiques
            \foreach \cx/\cy [count=\i] in {#5} {
                \pgfmathsetmacro{\currentcy}{{#6}[\i-1]}
                \pgfmathsetmacro{\diameter}{{#7}[\i-1]}
                \node[draw, dashed, circle, minimum size=\diameter] at (\cx cm, \currentcy cm) {};
            }
        \end{scope}

        % Points avec E_i dans un cercle bleu
        \foreach \x/\y [count=\i] in {#3} {
            \pgfmathsetmacro{\currenty}{{#4}[\i-1]}
            \node[draw=myblue, circle, inner sep=2pt, fill=white] at (\x cm, \currenty cm) {\scalebox{0.8}{\textcolor{myblue}{\small $\i$}}};
        }
    \end{scope}
    \draw[->, dashed, preaction={draw=white,line width=6pt}] (policy#2) -- ($(box#2.north)+(0,0)$);
}

\begin{figure}[t]
    \centering
    \resizebox{\textwidth}{!}{
    \begin{tikzpicture}[
        policybox/.style={
            draw, rectangle, text centered,
            minimum height=1cm, minimum width=2.5cm,
            align=center, rounded corners=0.2cm
        },
        pointset/.style={
            draw, minimum width=4cm, minimum height=4cm
        },
        circlepattern/.style={
            draw, dashed, circle, minimum width=1.5cm, minimum height=1.5cm
        }
    ]

    % --- Utilisation de la macro pour les trois policies ---
    \surrogatepolicy
    {-6cm}
    {1}
    {-1.5, -1.25, -1, 1.35, 1.0, -0.25, 0.25}
    {-0.75, -1.5, -1, 0.5, 0.0, 1.6, 1.6}
    {-1.25, 1.2, 0}
    {-1.1, 0.2, 1.75}
    {1.5cm, 1.25cm, 1.25cm}

    \surrogatepolicy
    {0cm}
    {2}
    {1, 1.3, 1.5, 0.1, 0.25, 0.45, 1.2}
    {1.5, 1, 1.5, 0.75, 0.2, -1.5, -1}
    {1.25, 0.1, 0.8}
    {1.3, 0.4, -1.2}
    {1.3cm, 1.2cm, 1.6cm}

    \surrogatepolicy
    {6cm}
    {3}
    {0.0, -0.5, 0.25, -1, -1.5, 1.6, 1.2}
    {0.0, 0.5, 0.5, -1.5, -1, 1.2, 1.6}
    {-1.25, -0.15, 1.5}
    {-1.25, 0.3, 1.5}
    {1.3cm, 1.5cm, 1.3cm}

    % --- Acolade et équation en dessous ---
    \draw[
        line width=1pt,  % Épaisseur du trait
        decorate,
        decoration={
            brace,
            amplitude=20pt,  % Hauteur de l'acolade
            mirror,          % Acolade sous la ligne
            raise=10pt       % Décalage vertical
        }
    ]
        (-8.5, -5.5) -- (8.5, -5.5);

    \node[below=0.75cm] at (0, -6) {
        \scalebox{1.5}{$\displaystyle
        \text{\ClusteringUniversalityName} = \frac{1}{\binom{n}{2}} \sum_{i=1}^{n} \sum_{j=0}^{i-1} \operatorname{ARI}(\mathcal{C}_{i}, \mathcal{C}_{j})
        $}
    };

    \end{tikzpicture}
    }
    \caption{Three point clouds obtained by projecting the same observations through three independently trained surrogate policies $\textcolor{myred}{\pi'_1}$, $\textcolor{myred}{\pi'_2}$, and $\textcolor{myred}{\pi'_3}$. The numbered circles 1 to 7 denote the same seven observations across the three point clouds, and the dashed contours denote the clusters recovered in each latent space.}
    \label{fig:clustering_universality_illustration}
\end{figure}
